\documentclass[10pt,a4paper]{amsart}
\usepackage[margin=19mm]{geometry}
\usepackage{amsmath,amssymb,mathtools}
\usepackage[scr=rsfs,cal=euler]{mathalfa}
\usepackage{xfrac}
\usepackage[unicode,hidelinks,bookmarks=false]{hyperref}
\newcommand{\ie}{i.e.\ }
\newcommand{\cf}{cf.\ }
\newcommand{\resp}{respectively\ }
\newcommand{\Subsection}{\subsection}
\providecommand{\nobreakdash}{\nobreak\hbox{-}}
\setlength{\emergencystretch}{2em}
\multlinegap0pt
\usepackage{amssymb}
\usepackage{enumerate}
\usepackage{tikz}
\usepackage[shortlabels]{enumitem}
\newtheorem{lem}{Lemma}
\newtheorem{prop}{Proposition}
\usepackage{cleveref}
\crefname{lem}{Lemma}{Lemmas}
\crefname{prop}{Proposition}{Propositions}
\newcommand{\cR}{{\mathcal R}}
\newcommand{\cV}{{\mathcal V}}
\newcommand{\pp}{\mathbb{P}}
\title{Adjoints of Polytopes: Determinantal Representations and Smoothness}
\author{Clemens Br\"user, Mario Kummer, Dmitrii Pavlov}
\date{Source: arXiv:2507.01672v1 (CC BY 4.0)}
\begin{document}
\maketitle

\noindent\emph{Source and licence.} Adjoints of Polytopes: Determinantal Representations and Smoothness. Version arXiv:2507.01672v1 (CC BY 4.0), distributed by arXiv under the Creative Commons Attribution 4.0 International licence, http://creativecommons.org/licenses/by/4.0/, as recorded in the arXiv licence metadata for this exact version and shown on the abstract page https://arxiv.org/abs/2507.01672v1. Full licence notice: \texttt{licenses/arxiv-2507.01672-CC-BY-4.0.txt}.\par\smallskip

\noindent\textit{Editorial scope.} Counted unit: the article's lem eq:adjpol\_expl (source0.tex lines 1047--1052). The article's complete original proof of it is delivered with it (source0.tex lines 1208--1281, one \texttt{proof} environment of 74 lines, which the article places away from the statement). No mathematics is written, repaired or re-derived: the statement and the proof are the article's own lines, in the article's own notation, and no retained line is substituted or re-flowed.  Retained uncounted as the source-local context the proof needs: lines 1057--1059; lines 1178--1180; lines 1201--1206.  One declared adaptation: exactly 1 ancillary strings inside the retained spans are omitted (image-inclusion commands and, where the source repeats a label, the repeated label definitions) and no other line differs from the source; the figure environments, with their placement, captions and labels, are kept, and the package records the omitted strings in fidelity receipts bound to the affected bodies.  No retained line contains a citation, so the wrapper carries no bibliography and no bibliography entry.  The retained lines contain the article's own diagram code, which the wrapper typesets with the standard tikz packages; no image file is delivered and the package declares no asset dependency.  Editorial formatting only, in addition: an AMS \texttt{amsart} preamble that restates the article's own declarations and macros used by the retained lines, namely the environments \texttt{lem}, \texttt{prop} on separate counters, together with the standard packages those lines need.  The retained environments print in this document as Lemma 1, Proposition 1 in order of appearance, whereas the article prints the same items with the numbering of its own full document.  The complete native source of the article as distributed in the arXiv e-print for this exact version is bundled unchanged as \texttt{sources/180897/source0.tex}.
\begin{lem} \label{lem:AdjTangent}
    Let $i,j\in\{1,\ldots,n\}$ such that $\{q\} = L_i \cap L_j\subseteq\cR(P)$, and let $Q$ denote the quadrilateral spanned by the vertices $v_{i-1}, v_i, v_{j-1}, v_j$. Then $\cV(\alpha_{Q})$ is the tangent line to $\cV(\alpha_P)$ at $q$.
\end{lem}

\begin{prop}\label{thm:adjointcontact}
    Let $P$ be a convex polygon with $n$ vertices and adjoint curve $A = \cV(\alpha_P)$. Let $P'$ be the polygon obtained from $P$ by removing the vertex $v_n$. The intersection of the adjoint curve $A'$ of $P'$ with $A$ is equal to $\cR(P) \smallsetminus (L_1 \cup L_n)$. Moreover, each intersection point is of multiplicity two.
\end{prop}

\begin{figure}[h]
    \centering
    
    \caption{The blue line is the adjoint (quartic) curve of the heptagon. The red and black lines are the adjont curves of the correspondingly colored subpolygons.}
    \label{fig:WLOGsubquadril}
\end{figure}

\begin{lem}
    In the notation above, the adjoint polynomial of the convex polygon $P$ is given by
    \begin{equation} \label{eq:adjpol_expl}
        \alpha_P(x) = \sum_{i = 1}^n \det(w_i, w_{i+1}) \prod_{j \not\in\{ i, i+1\}} l_j.
    \end{equation}
\end{lem}

\begin{proof}
 After choosing a suitable affine chart in $\pp^2$ and applying a linear coordinate change we may without loss of generality assume that
 \begin{equation*}
     v_{n-2} = (-\gamma,\gamma), v_{n-1} = (-1,1), v_n = (1,1), v_1 = (\gamma,\gamma)
 \end{equation*}
 for some $\gamma > 1$ (see \Cref{fig:WLOGsubquadril}). Then $A_Q = \cV(x_2)$ is the $x_1$-axis. Writing $\alpha_k$ for the adjoint polynomial of $A_k$ as defined in \Cref{eq:adjpol_expl}, it suffices to prove that
 \begin{equation*}
     -\gamma\alpha_n(x_1,0) = x_1^2\alpha_{n-2}(x_1,0).
 \end{equation*}
 Recall that $l_i(x_1,x_2) = \langle w_i, \binom{x_1}{x_2} \rangle + c_i \geq 0$ are the facet inequalities of $P$. For the sake of readability we write $w_i = (a_i,b_i)^t$. We thus have 
 \begin{equation*}
 \begin{split}
     \alpha_n(x_1,0) &= \sum_{i=2}^{n-3} \det(w_i,w_{i+1}) \left(\prod_{2\leq j\leq n-2, j\neq i,i+1} l_j(x_1,0)\right) x_1^2 \\
     &+ (a_{n-2}-b_{n-2})\left(\prod_{2\leq j \leq n-3} l_j(x_1,0)\right) x_1\\ & + (a_2+b_2)\left(\prod_{3\leq j \leq n-2} l_j(x_1,0)\right) x_1\\
     &- \left(\prod_{2\leq j \leq n-2} l_j(x_1,0)\right) x_1 + \left(\prod_{2\leq j \leq n-2} l_j(x_1,0)\right) x_1 \\
    &= x_1^2 \cdot(\dots)\\ & + x_1 \left(\prod_{3\leq j\leq n-3} l_j(x_1,0)\right)((a_{n-2}-b_{n-2})l_2(x_1,0) + (a_2+b_2)l_{n-2}(x_1,0)).
 \end{split}
 \end{equation*}
 Similarly, we have
 \begin{equation*}
 \begin{split}
     \alpha_{n-2}(x_1,0) &= -\gamma \sum_{i=2}^{n-3} \det(w_i,w_{i+1}) \prod_{2\leq j\leq n-2, j\neq i,i+1} l_j(x_1,0) \\
     & + a_{n-2} \prod_{2\leq j\leq n-3} l_j(x_1,0) - a_2 \prod_{3\leq j\leq n-2} l_j(x_1,0) \\
     &= -\gamma(\dots) + \left(\prod_{3\leq j\leq n-3} l_j(x_1,0)\right)(a_{n-2}l_2(x_1,0) - a_2l_{n-2}(x_1,0)).
 \end{split}
 \end{equation*}
 It now suffices to prove that
 \begin{equation*}
     -\gamma((a_{n-2}-b_{n-2})l_2(x_1,0) + (a_2+b_2)l_{n-2}(x_1,0)) = x_1(a_{n-2}l_2(x_1,0) - a_2l_{n-2}(x_1,0)).
 \end{equation*}
 Using $-\gamma a_{n-2} + \gamma b_{n-2} + c_{n-2} = 0$ and $\gamma a_2 + \gamma b_2 + c_2 = 0$ we may rewrite all occurrences of $c_{n-2}$ and $c_2$ (which are implicit in $l_{n-2}$ and $l_2$) in terms of $a_{n-2},a_2,b_{n-2},b_2$, which yields
 \begin{equation*}
 \begin{split}
     &(a_{n-2} - b_{n-2})l_2(x_1,0) + (a_2 + b_2)l_{n-2}(x_1,0)  \\ =&
     (a_{n-2} - b_{n-2})(a_2x_1-\gamma(a_2+b_2)) + (a_2 + b_2)(a_{n-2}x_1 + \gamma(a_{n-2}-b_{n-2}))  \\ =&
     (2a_2a_{n-2} - a_2b_{n-2} + a_{n-2}b_2)x_1
 \end{split}
 \end{equation*}
 and
 \begin{equation*}
 \begin{split}
     &a_{n-2}l_2(x_1,0) - a_2l_{n-2}(x_1,0)  \\ =
     &a_{n-2}(a_2x_1 - \gamma(a_2+b_2)) - a_2(a_{n-2}x_1 + \gamma(a_{n-2} - b_{n-2}))  \\ =&
     - \gamma(2a_2a_{n-2} - a_2b_{n-2} + a_{n-2}b_2).
 \end{split}
 \end{equation*}
 Thus the first claim follows.

 Now let $A_n$ be smooth. For proving that  $A_{n-2} \cup A_Q$ is a contact curve to $A_n$, by B\'{e}zout's theorem it suffices to count points of contact with $A_n$. From \Cref{lem:AdjTangent} we know that every subquadrilateral of $P_{n-2}$ that shares two non-adjacent edges with both $P_n$ and $P_{n-2}$ gives one such point of contact. Every such quadrilateral is uniquely determined by two such edges of $P_{n-2}$. This gives rise to exactly
 \begin{equation*}
     \binom{n-3}{2} - \underbrace{(n-4)}_\textrm{unordered pairs of adjacent edges}
 \end{equation*}
 many such quadrilaterals and thus points of contact.
 Furthermore, the line $A_Q$ is tangent to $A_n$ at one point, and intersects $A_n$ in $n-5$ further points that all lie on $A_{n-2}$ (as shown in the first part of the proof). In total this gives rise to
 \begin{equation*}
     \binom{n-3}{2} - (n-4) + (n-5) + 1 = \frac{(n-3)(n-4)}{2}
 \end{equation*}
 contact points and we conclude by B\'{e}zout's theorem that $A_{n-2} \cup A_Q$ is a contact curve to $A_n$. Its contact divisor can be determined from the same considerations. 
 Namely, we have seen above that
 \begin{align*}
     A_Q.A_n&=2T_{1,n-1}+(R_1+\cdots+R_{n-5}),\textrm{ and}\\
     A_{n-2}.A_n&=2\sum_{i=2}^{n-4}\sum_{j=i+2}^{n-2} T_{ij}+(R_1+\cdots+R_{n-5}),
 \end{align*}
 where $T_{ij}$ is the intersection point of $L_i$ and $L_j$ and $T_{1,n-1},R_1,\ldots,R_{n-5}$ are the intersection points of $A_Q$ and $A_n$. Furthermore, we have
 \begin{align*}
     L_{n-1}.A_n&=T_{1,n-1}+\sum_{i=2}^{n-3}T_{i,n-1},\textrm{ and}\\
     A_{n-1}.A_n&=2\sum_{i=2}^{n-3}\sum_{j=i+2}^{n-1}T_{ij},
 \end{align*}
 where the second equality follows from \Cref{thm:adjointcontact}. From this we obtain
 \begin{equation*}
     A_{n-1}.A_n+A_Q.A_n=A_{n-2}.A_n+2L_{n-1}.A_n
 \end{equation*}
 which is equivalent to the claim.
\end{proof}

\end{document}
